\scanpage{5}
\mspage{3}
i.e. \quad a) $C^\bullet_{\mathrm{DRS}}(X,\mathbb Q)$ dépend de $X$ via $S_{\unsure{\ast}}(X)$

b) Sur $(\mathrm{Ss})\deleted{\unsure{\ast}}=(\mathrm{Ens})_{\ast}$, le foncteur $C^\bullet_{\mathrm{DRS}}(-,\mathbb Q)$ est représentable par $\mathrm{DRS}^{\bullet}_{\ast}$.

Or \quad a) $\mathrm{DRS}^{\bullet}_{\ast}$ est une résolution de $\mathbb Q_{\ast}$ \inserted{(dans la cat. des groupes $1/2$ simpl.)} (lemme de Poincaré algébrique sur les espaces \inserted{$\mathbb Q$-}affines $E^{[n]}$)

b) Les composantes $\mathrm{DRS}^{i}_{\ast}$ $(i\geq 0)$ de $\mathrm{DRS}^{\bullet}_{\ast}$ sont des objets \inserted{\unsure{ab}} acycliques du topos $\mathrm{Ss}$ (ce qui revient à dire que leurs $\pi_j(\mathrm{DRS}^{i}_{\ast})$ sont nuls, ce qui se vérifie facilement. \inserted{Cela serait itou avec des formes $C^\infty$ à coeff. ds $\mathbb R$ ou $\mathbb C$, ou analytiques réelles, \deleted{ou} analytiques complexes, \ldots)}

\deleted{\unsure{On peut partir du}} On aimerait avoir une $\mathbb Z$-algèbre diff. $C_{\mathrm{DR}}(X,\mathbb Z)$ \uline{alternée} (dépendant fonct. de $X$, ou $S_{\ast}(X)$) dont la coh. soit la coh. entière de $\mathbb Z$. Mais il paraît qu'il y a des ths contre ! On va pourtant y arriver \deleted{\unsure{quand}}\unsure{en un sens|en q. sens} !

Si on prend $C_{\mathbb Z\text{-}\mathrm{DR}}^{\bullet}(E^{[n]})$ (car $E^{[n]}$ a une $\mathbb Z$-structure affine) c'est a) qui devient déjà faux : pour intégrer $\int x^n dx$, il faut un dénominateur \inserted{avec} $\frac{x^{n+1}}{n+1}$ ! Mais en géom. alg. on est déjà familiarisé avec une façon de sauter à pieds-joints \deleted{sur}\inserted{par dessus} le caniveau, en introduisant des p.d. et des polynômes \inserted{(ou séries formelles)} à puiss. div. Si $E^{[n]}$ avait une origine sur $\mathbb Z$ (i.e. provenait can. d'un $\mathbb Z$-module libre de déf.) on aurait un complexe de De Rham à p.d. Mais pas pour un espace affine !. Notons
\[
\begin{aligned}
\mathrm{DRS}^{\bullet}_{[n]}&\simeq\deleted{\mathbb Q[X_i,dX_i]/}\ C^\bullet_{\mathrm{DR}/\mathbb Q}\big(\mathbb Q[(X_i)_{0\leq i\leq n}]/\textstyle\sum X_i-1\big)\\
&\simeq\mathbb Q[X_i,dX_i]_{0\leq i\leq n}/\big(\textstyle\sum X_i-1,\ \textstyle\sum dX_i\big)\\
&\simeq\mathbb Q[X_0,\ldots,X_{n-1}][dX_0,\ldots,dX_{n-1}]
\end{aligned}
\]

% archive sheet number: 4
